\cleardoublepage
\chapter*{附录 A\quad 外文文献翻译}
\addcontentsline{toc}{chapter}{附录 A\quad 外文文献翻译}

\begin{center}
{\heiti\fontsize{14pt}{20pt}\selectfont 基于物理引导的多变量图建模与故障传播分析的风力发电机组根因定位}
\end{center}

\vspace{6pt}
\begin{center}
冯晨龙$^{a}$，刘超$^{a,b,*}$，蒋东翔$^{a,c}$

\vspace{4pt}
{\small
$^{a}$ 清华大学能源与动力工程系，北京 100084 \\
$^{b}$ 教育部热科学与动力工程重点实验室，清华大学，北京 100084 \\
$^{c}$ 电力系统及发电设备控制和仿真国家重点实验室，清华大学，北京 100084 \\
}
\end{center}

\vspace{6pt}
\noindent {\small \textbf{原文信息：}Feng C, Liu C, Jiang D. Root cause localization for wind turbines using physics guided multivariate graphical modeling and fault propagation analysis[J]. Knowledge-Based Systems, 2024, 295: 111838.}

\vspace{12pt}

\section*{A.1\quad 摘要}
\addcontentsline{toc}{section}{A.1\quad 摘要}

在风能广泛利用的背景下，发电效率的提高和运维成本的降低增加了通过故障传播分析对风力发电机组进行异常检测和定位的需求，而SCADA系统采集的大规模数据为此提供了充足的数据支撑和有力的帮助。基于多变量监测数据构建正常运行工况下的高精度风力发电机组数字孪生模型，并基于预测偏差和回溯分析实现根因定位，是一种有效且经济的解决方案。本文提出了一种基于显式-隐式知识融合和多变量时间序列图神经网络的风力发电机组根因定位框架，充分利用图结构学习对变量间依赖关系的挖掘能力和时间序列图神经网络的时序预测优势。基于显式规则构建了五类显式知识以描述监测变量之间的显式关系。通过基于图结构学习思想的准硬注意力机制，将显式知识与数据中隐含的隐式知识进行融合，获得融合知识（即图邻接矩阵）。基于融合知识，利用多变量时间序列图神经网络构建了高精度风力发电机组数字孪生模型。基于预测偏差定义了监测变量状态的异常程度，据此可准确描述故障状态随时间的传播过程，并通过回溯分析实现根因定位。利用某风电场的实际故障数据、某糖厂和某火力发电厂的数据进行了案例研究。结果表明，所提出的框架能够为研究故障传播过程和根因定位提供解决方案，并展现出一定的鲁棒性和泛化能力。

\textbf{关键词：}风力发电机组；数据采集与监控系统；多变量时间序列预测；图神经网络

\section*{A.2\quad 引言}
\addcontentsline{toc}{section}{A.2\quad 引言}

凭借清洁低碳的优势，风力发电在应对气候变化和环境污染方面做出了卓越贡献。2022年，全球新增风电装机容量达到77.6~GW，累计装机容量为906~GW。根据《2023全球风能报告》，预计到2024年全球陆上风电新增装机容量将首次突破100~GW，到2025年海上风电新增装机容量将达到25~GW的新高。随着风电的快速发展，对风力发电机组进行有效状态监测和异常检测的需求日益增加。近年来，大多数风电场都配备了数据采集与监控（SCADA）系统，其采集的数据具有规模大、成本低、监测覆盖面广等优势，使其成为实施风力发电机组智能运维的最经济选择。然而，该系统目前仅用于状态监测和报警，未能使存储的大量历史数据创造经济价值。此外，现有的许多研究侧重于基于SCADA数据对风力发电机组进行异常检测，但难以识别根因所在位置。因此，工业界和学术界都迫切希望基于SCADA数据实现风力发电机组的根因定位，以降低运维成本并提高风电场的经济效益。

考虑到SCADA数据采样率较低（通常为10分钟）且相比振动数据包含的信息量有限，不便于采用类似处理振动数据的方法（如傅里叶变换）来处理该数据。在大多数情况下，SCADA数据描述的是设备的正常状态，异常状态所占比例很小。因此，可以考虑基于多变量监测数据构建描述设备正常运行工况的高精度数字孪生模型，然后通过预测偏差和回溯分析实现异常状态的根因定位。该方案的挑战在于如何以更好反映实际过程的方式构建描述设备正常状态的高精度数字孪生模型。"更好反映实际过程"意味着由各种传感器监测的复杂系统可以被抽象为监测变量和变量间依赖关系。换言之，一个复杂的机械系统可以由监测变量（节点）及其关系（边）组成的复杂图结构来描述，每个节点状态的更新依赖于捕获历史时间序列信息和探索相邻变量的空间结构信息。因此，该数字孪生模型的构建可以在学术上概括为基于图的多变量时间序列预测模型的构建，该模型能够以更好反映实际过程的方式将多变量时间序列数据抽象为变量内时序依赖关系和变量间空间依赖关系的组合，并对所有变量给出预测。

为实现高精度多变量时间序列预测，许多基于典型深度学习方法的时间序列预测模型做出了努力。然而，这些模型主要考虑时序依赖关系对预测结果的影响，对变量间交互（即变量间空间依赖关系）关注较少或仅以隐式方式关注，导致模型获取的信息不完整。图结构数据以其显式的节点连接可以自然地解决这一问题。因此，一方面需要将多变量时间序列数据抽象为图结构数据，其中监测变量作为节点，变量间的隐式相互关系作为边；另一方面，需要利用图神经网络在处理图结构数据方面的出色能力来构建多变量时间序列图神经网络（MTGNN），以进一步提高预测精度。

将多变量监测数据抽象为图可以从两个角度实现：规则驱动（或经验驱动）和数据驱动。规则驱动方法完全依赖物理规则或专家经验来描述变量间依赖关系以构建图，但其难以实施，且获得的知识通常不能完全准确地表达变量之间的关系。数据驱动方法是从数据中学习变量间依赖关系，也称为图结构学习（GSL），原因是大多数场景不能满足"使用预定义图进行模型训练"的要求。由于大量历史监测数据中隐含的信息遵循稳定的统计模式，因此该方法获得的隐式知识可以部分正确地表达变量之间的依赖关系。因此，可以考虑在从数据获取隐式知识的过程中自适应地融入不同的显式知识，以基于物理和数据联合学习更准确的图结构。

具体而言，为实现上述构建高精度数字孪生模型和实现根因定位的方案，本文结合基于物理和数据的GSL以及MTGNN时间序列预测的优势，提出了一种基于显式-隐式知识融合和MTGNN的风力发电机组根因定位框架。首先，设计了基于不同规则的显式知识，从物理和数据的角度充分探索变量间依赖关系，并通过准硬注意力机制将风力发电机组的显式和隐式知识进行融合，以更准确地将风力发电机组抽象为变量及其交互（即图）。其次，构建了高精度数字孪生模型，即基于多变量时间序列数据的MTGNN模型，以准确及时地描述风力发电机组的正常运行状态，并基于该模型设计了监测变量状态的异常程度。最后，获得了故障状态在不同变量间随时间的传播（即故障传播分析），并通过反向时间轴上的回溯分析找到故障的根节点以实现根因定位。

本文的主要贡献如下：

\begin{enumerate}
\item[i)] 设计了基于不同规则的显式知识，包括机械传动关系、温度影响关系、相关性、互信息和因果性，用于描述风力发电机组的变量间依赖关系，并使用准硬注意力机制在GSL框架中融合显式和隐式知识，获得用于风力发电机组的准确的物理与数据融合图。

\item[ii)] 利用物理与数据融合图和MTGNN构建高精度风力发电机组数字孪生模型，使用设计的监测变量状态异常程度基于预测偏差进行故障传播分析，最终通过回溯分析实现根因定位。

\item[iii)] 利用实际风电场的案例验证了所提框架的有效性，并利用实际糖厂和火力发电厂在不同应用场景下的案例进一步验证了该框架的鲁棒性和泛化能力。
\end{enumerate}

本文其余部分组织如下。第二节概述了GSL、多变量时间序列预测和根因定位的研究进展。第三节介绍了包括因果性和MTGNN在内的预备知识。第四节提出了基于显式-隐式知识融合和MTGNN的风力发电机组根因定位框架。第五节描述了案例研究所需的数据集、实验设置和评价指标，并对结果进行了展示和讨论。最后在第六节总结了结论。

\section*{A.3\quad 相关工作}
\addcontentsline{toc}{section}{A.3\quad 相关工作}

\subsection*{A.3.1\quad 图结构学习（GSL）}

图结构学习旨在联合学习优化的图结构及其相应的图表示。典型的GSL模型包含两个可训练部分：一是GNN编码器，用于接收图作为输入并生成下游任务所需的嵌入向量；二是结构学习器，用于建模图中节点之间的连接。一般来说，GSL的流程包括三个阶段：图构建、图结构建模和消息传播。其中，图结构建模是GSL的核心，可以建模边连接以细化初始图，并可分为三类：基于度量的方法、基于神经网络的方法和直接方法。

基于度量的方法使用不同的核函数计算节点特征/嵌入的相似度作为边权重。这些方法主要基于网络同质性假设，即边倾向于连接相似节点，并通过促进类内连通性来细化图结构以产生更紧致的表示。AGCN是基于距离度量学习的GSL模型，使用广义马氏距离计算相似度，然后使用高斯核函数细化初始图。GRCN、GAUG-M和CAGNN使用两个节点嵌入的内积（内积核）来建模边权重。类似地，余弦核、扩散核和多核融合函数也可用于边权重建模。

基于神经网络的方法使用更复杂的神经网络来建模边权重。例如，单层神经网络和多层感知器等简单神经网络可用于GSL。随着近年来注意力机制的快速发展，一些方法如GaAN、PGN和MAGNA利用不同的注意力机制动态放大或衰减连接以生成邻接矩阵。相比之下，基于Transformer的方法（如Graphomer和GraphiT）不仅考虑一阶邻接，而是将所有节点视为邻居，并将图结构信息编码到图Transformer中，然后用注意力加权边。

与上述方法不同，直接方法不再依赖节点表示进行边权重建模，而是将邻接矩阵视为GNN训练中的可学习自由变量，这提供了更大的灵活性，但在学习矩阵参数方面提出了更大的困难。具体而言，一些研究使用正则化器来优化邻接矩阵，而另一些则使用概率方法来构建邻接矩阵，即假设内在图结构是从某种分布中采样的。例如，LDS-GLNN通过从伯努利分布中采样来构建邻接矩阵。

在本文中，受正则化图结构学习（RGSL）的启发，使用内积核度量节点嵌入之间的距离来构建初始邻接矩阵，然后通过Gumbel Softmax对初始图结构进行离散采样以生成所需的隐式图结构。

\subsection*{A.3.2\quad 基于GNN的多变量时间序列预测}

传感和数据流处理技术的进步催生了时间序列数据的爆发式增长，使其成为许多领域中最常见的数据类型之一。时间序列数据不仅可用于关注过去的趋势，还可用于预测、分类、异常检测和数据插补等多项任务。时间序列预测作为时间序列分析的重要课题之一，旨在基于历史观测值预测未来值。传统的时间序列预测主要基于统计回归模型，如自回归积分滑动平均（ARIMA）、向量自回归（VAR）和向量自回归滑动平均（VARMA）等，将预测的未来值视为历史值的线性组合。这类方法由于"线性组合"假设的天然缺陷，在面对复杂的时间序列关系（如非线性和多变量序列间关系）时难以达到令人满意的预测性能。近年来，基于深度学习的方法，包括RNN、CNN和Transformer，凭借其卓越的非线性时空特征捕获能力，在时间序列预测方面取得了巨大成功。

然而，考虑到多变量时间序列数据涉及时间和变量之间的复杂交互，上述深度学习方法最大的局限之一是无法显式建模非欧几里得空间中时间序列之间存在的空间关系，这限制了其表达性能。因此，有效学习这些交互复杂性的解决方案是使用由网络或图构建的模型，其中变量作为节点，变量间的复杂关系作为边。在此背景下，GNN作为学习非欧几里得数据表示的强大工具，在显式有效地建模多变量时间序列中的时空依赖关系方面展现出巨大潜力，从而增强预测性能。

基于GNN的时间序列预测模型在方法论上可分为两类：空间依赖关系建模和时序依赖关系建模。

空间依赖关系，即序列间关系，对预测性能有显著影响。当前研究从三个角度建模空间依赖关系：谱域GNN、空间域GNN和混合模型。基于谱域GNN的模型主要利用切比雪夫多项式近似图卷积来建模序列间依赖关系，如STGCN和StemGNN模型。基于空间域GNN的模型使用消息传递或图扩散来建模依赖关系。例如，DCRNN和GraphWaveNet通过将图扩散层结合到GRU或时序卷积中来建模时间序列数据。后续工作包括STSGCN、STFGNN、MTGNN和ASTTGN等，从不同角度创新性地建模时空依赖关系以增强学习能力。通过融合谱域GNN和空间域GNN构建混合模型是另一条创新路线。

时序依赖关系建模是多变量时间序列预测的另一个重要元素，其主要目标是学习一个有效的时序模型以准确捕获时间序列中数据点之间的依赖关系。基于循环的模型、基于卷积的模型、基于注意力的模型和混合模型均可用于此目的。基于循环的模型用于在一些早期方法中理解时间间依赖关系，如DCRNN将图扩散与GRU结合来建模时空依赖关系。基于卷积的模型为时间间依赖关系建模提供了更有效的解决方案，例如STGCN采用整合1-D卷积和门控线性单元的时序门控卷积。MTGNN利用多种卷积核尺寸来增强时序卷积。近来，越来越多的研究关注基于注意力的模型。混合模型是建模时间间依赖关系的另一种应用。

\subsection*{A.3.3\quad 根因定位}

目前，有许多基于SCADA数据实现风力发电机组异常检测的方法，但其中很少关注根因定位。基于多变量时间序列信息实现风力发电机组的根因定位是一项重要且具有挑战性的任务。然而，现有的许多根因定位研究集中在电站电路、模拟电路和可靠性评估等方面，针对风力发电机组的研究相对较少。多变量时间序列图神经网络的发展为根因定位提供了新思路，即在基于多变量时间序列图神经网络的正常工况下建立高精度模型，当故障状态在系统中传播时，各监测变量偏离正常状态的先后顺序不一致，因此当通过回溯分析明确这一顺序时即可实现根因定位。

\section*{A.4\quad 预备知识}
\addcontentsline{toc}{section}{A.4\quad 预备知识}

\subsection*{A.4.1\quad 因果性}

在许多实际应用中，从监测数据中识别因果结构是一项重要但具有挑战性的任务。Wang等人提出了一种名为基于因果发现排序的强化学习（CORL）的因果结构发现算法，利用变量排序的缩减空间和强化学习强大的搜索能力，并验证了该算法优于现有的因果发现算法，因此本研究使用该方法从风力发电机组的监测数据中提取因果结构。

CORL的架构由编码器模块和解码器模块组成，并使用强化学习进行神经网络优化（如图~A-1所示）。

\textbf{编码器：}用于将监测数据映射到嵌入状态空间。

\textbf{解码器：}用于将状态空间映射到动作空间。

\textbf{优化：}CORL的优化目标是学习一个最大化期望累积增益的策略。

\textbf{变量选择：}对于上述方法获得的变量排序，CORL利用不同的变量选择方法获得最终的因果结构图。

% 图A-1: CORL模型的策略说明
% \begin{figure}[htbp]
% \centering
% \includegraphics[width=0.7\textwidth]{figures/appendixA/fig1_CORL.png}
% \caption{CORL模型的策略说明}
% \label{fig:appA-CORL}
% \end{figure}

\subsection*{A.4.2\quad 多变量时间序列图神经网络（MTGNN）}

MTGNN是一种专为多变量时间序列数据设计的通用图神经网络框架，由图学习模块、图卷积模块、时序卷积模块和跳跃连接与输出模块组成。然而，在本文中，图学习模块被去除，改为使用从显式和隐式知识中学习到的预定义图。本文使用的MTGNN框架如图~A-2所示。首先，使用$1 \times 1$标准卷积将输入投影到潜在空间，然后通过时序卷积模块和图卷积模块交替捕获时序和空间依赖关系。残差连接和跳跃连接用于避免梯度消失。最后，隐藏特征由输出模块按照所需维度投影以产生结果。

% 图A-2: 本文使用的MTGNN框架
% \begin{figure}[htbp]
% \centering
% \includegraphics[width=0.9\textwidth]{figures/appendixA/fig2_MTGNN.png}
% \caption{本文使用的MTGNN框架}
% \label{fig:appA-MTGNN}
% \end{figure}

\subsubsection*{A.4.2.1\quad 图卷积模块}

图卷积模块可以融合节点信息与其邻居信息以捕获空间依赖关系。该模块由两个混合跳跃传播层组成，通过信息传播步骤和信息选择步骤处理流经每个节点的信息流。

具体而言，给定图结构$A$，信息传播步骤将递归地传递前一步的节点隐藏状态$\boldsymbol{H}^{(k-1)}$以及原始状态$\boldsymbol{H}_{in}$，以获得处理后的隐藏状态$\boldsymbol{H}^{(k)}$。设计了超参数$\beta$来控制保留根节点原始状态的比例，以避免多层图卷积堆叠导致的隐藏状态收敛问题。数学上，信息传播步骤定义为：
\begin{equation}
\boldsymbol{H}^{(k)} = \beta \boldsymbol{H}_{in} + (1-\beta)\tilde{\boldsymbol{A}}\boldsymbol{H}^{(k-1)}
\end{equation}
其中$\tilde{\boldsymbol{A}} = \tilde{\boldsymbol{D}}^{-1}(\boldsymbol{A}+\boldsymbol{I})$，$\tilde{\boldsymbol{D}}_{ii} = 1 + \sum_j A_{ij}$，$\boldsymbol{I}$为单位矩阵，$\boldsymbol{D}$为度矩阵，$k \in [1, K]$为当前层，$K$为传播深度。

信息选择步骤使用可学习的信息注意力矩阵$\boldsymbol{W}^{(k)}$对隐藏状态信息进行过滤以保留最重要的信息，定义为：
\begin{equation}
\boldsymbol{H}_{out} = \sum_{i=0}^{K} \boldsymbol{H}^{(k)}\boldsymbol{W}^{(k)}
\end{equation}
其中$\boldsymbol{H}_{out}$表示当前层的输出隐藏状态，$\boldsymbol{H}^{(0)} = \boldsymbol{H}_{in}$。

\subsubsection*{A.4.2.2\quad 时序卷积模块}

时序卷积模块可以通过一组标准膨胀1D卷积滤波器提取高级时序特征。该模块由两个膨胀inception层组成，其中一个后接tanh激活函数用作滤波器，另一个后接sigmoid激活函数用作控制传递到下一模块的信息量的门。

对于使用1D卷积滤波器提取序列模式的时序卷积模块，多尺寸滤波器有助于发现各种尺度的时域模式，膨胀卷积层有助于降低处理长序列的模型复杂度。因此，引入了$1 \times 2$、$1 \times 3$、$1 \times 6$和$1 \times 7$四种尺寸的滤波器构成时序inception层，并提出了结合inception和膨胀的膨胀inception层。给定一维序列$\boldsymbol{z} \in \mathbb{R}^T$及由$f_{1\times2} \in \mathbb{R}^2$、$f_{1\times3} \in \mathbb{R}^3$、$f_{1\times6} \in \mathbb{R}^6$和$f_{1\times7} \in \mathbb{R}^7$组成的滤波器，膨胀inception层定义为：
\begin{equation}
\boldsymbol{z} = \mathrm{concat}(\boldsymbol{z} \star f_{1\times2},\ \boldsymbol{z} \star f_{1\times3},\ \boldsymbol{z} \star f_{1\times6},\ \boldsymbol{z} \star f_{1\times7})
\end{equation}
其中$\boldsymbol{z} \star f_{1\times k}$表示膨胀卷积，其定义为：
\begin{equation}
\boldsymbol{z} \star f_{1\times k}(t) = \sum_{s=0}^{k-1} f_{1\times k}(s) \boldsymbol{z}(t - d \times s)
\end{equation}
其中$d$为膨胀因子。

\subsubsection*{A.4.2.3\quad 跳跃连接与输出模块}

跳跃连接层是标准卷积：
\begin{equation}
\boldsymbol{h} \star f_{1\times L_i}(t) = \sum_{s=0}^{L_i - 1} f_{1\times L_i}(s) \boldsymbol{h}(t - s)
\end{equation}
其中$\boldsymbol{h}$表示该层的输入，$L_i$为第$i$个跳跃连接层的输入序列长度。

输出模块由两个$1 \times 1$标准卷积层组成，可将输入通道维度转换为所需的输出维度。在本文中，预测值为下一时步的值，因此所需的输出维度为1。

\section*{A.5\quad 方法论}
\addcontentsline{toc}{section}{A.5\quad 方法论}

本文提出了一种基于显式-隐式知识融合和多变量时间序列图神经网络的三阶段风力发电机组根因定位框架，如图~A-3所示。在第一阶段，分别定义了基于机械传动关系、温度影响关系、相关性、互信息和因果性的五类显式知识，并在GSL框架中使用准硬注意力机制将从数据中学习到的隐式知识与这些显式知识进行融合，以获得显式-隐式融合知识（即基于多变量监测数据的邻接矩阵）。在第二阶段，基于MTGNN利用融合知识构建风力发电机组的高精度数字孪生模型，并设计了描述监测变量状态的异常程度。在第三阶段，利用该异常程度获得风力发电机组在每个时刻的异常程度分布图，即故障状态随时间的传播，并通过回溯分析实现根因定位。

% 图A-3: 所提出的根因定位框架的总体架构
% \begin{figure}[htbp]
% \centering
% \includegraphics[width=\textwidth]{figures/appendixA/fig3_framework.png}
% \caption{所提出的根因定位框架的总体架构}
% \label{fig:appA-framework}
% \end{figure}

\subsection*{A.5.1\quad 问题建模}

在本文中，风力发电机组数字孪生模型的构建可以被抽象为基于MTGNN的多变量时间序列预测模型的构建。令$\boldsymbol{z}_t \in \mathbb{R}^N$表示$N$维变量在$t$时刻的值，其中$\boldsymbol{z}_t[i] \in \mathbb{R}$表示第$i$个变量在$t$时刻的值。给定该多变量变量长度为$P$的观测序列$\boldsymbol{X} = \{\boldsymbol{z}_{t_1}, \boldsymbol{z}_{t_2}, \ldots, \boldsymbol{z}_{t_P}\}$，预测目标为下一时刻的值$\boldsymbol{Y} = \{\boldsymbol{z}_{t_{P+1}}\}$。模型训练的目标是通过最小化绝对损失构建从$\boldsymbol{X}$到$\boldsymbol{Y}$的映射函数$f(\cdot)$。

本文中构建MTGNN模型所需的图相关定义如下：

\textbf{定义4.1}（图）：图在数学上表示为$G = (V, E)$，其中$V$为节点集合，$E$为边集合。本文中，节点为风力发电机组SCADA系统的监测变量，边为不同变量之间的交互。$N$为图中节点数量，即监测变量的数量。

\textbf{定义4.2}（节点邻域）：令$v \in V$表示一个节点，$e = (u, v) \in E$表示从$u$指向$v$的边，则节点$v$的邻域定义为$\mathcal{N}(v) = \{u \in V | (u, v) \in E\}$。

\textbf{定义4.3}（邻接矩阵）：邻接矩阵是图的数学表示，记为$\boldsymbol{A} \in \mathbb{R}^{N \times N}$，若$e_{ij} \in E$则$A_{ij} = 1$，否则$A_{ij} = 0$。

\subsection*{A.5.2\quad 物理引导的图建模}

图建模的主要目标是通过尽可能充分地正确提取不同监测变量之间的关系，将风力发电机组抽象为图。在本文中，物理引导的图建模主要从两个角度实现：显式知识和隐式知识。显式知识从五个方面描述，包括机械传动关系、温度影响关系、相关性、互信息和因果性。隐式知识的获取依赖于在GSL框架下学习数据中隐含的关系。此外，在获取隐式知识的过程中，使用准硬注意力机制融合显式知识和隐式知识，旨在以学习隐式知识为核心、以不同权重的显式知识为补充。

\subsubsection*{A.5.2.1\quad 显式知识}

显式知识是由显式规则定义的边。本文从五个方面定义了变量之间的关系（即边），包括机械传动关系、温度影响关系、相关性、互信息和因果性。

\textbf{i) 由机械传动关系定义的边}

这类边主要描述风力发电机组在将风能转化为电能的过程中不同子部件之间的能量传递关系。根据风力发电机组的结构，传动链中的能量传递是单向过程，因此这类边是单向的。具体而言，这种单向关系可以是一对多或多对一。

\textbf{ii) 由温度影响关系定义的边}

这类边主要描述环境温度对风力发电机组内部温度的影响以及机组内部不同位置温度之间的相互影响。值得注意的是，环境作为恒温无限热源，不会受到机组内部温度的影响，因此环境温度对内部温度的影响是单向的。相比之下，不同位置的温度之间的影响是双向的。此外，不同位置子部件温度之间的影响通过机舱温度联系，即彼此之间没有直接影响。

\textbf{iii) 由相关性定义的边}

这类边使用Pearson相关系数描述风力发电机组不同变量之间的线性相关性，选择阈值为0.7。

\textbf{iv) 由互信息定义的边}

这类边使用互信息描述风力发电机组不同变量之间基于信息论的统计依赖关系，选择阈值为0.7。

\textbf{v) 由因果性定义的边}

这类边使用CORL算法提取风力发电机组不同变量之间的因果强度。在学习过程中，使用窗口大小为三个月、步长为一个月的滑动窗口选择数据，用于学习不同变量之间的因果关系，最终通过对所有窗口结果取平均得到所需的因果图，选择阈值为0.7。

\subsubsection*{A.5.2.2\quad 隐式知识}

隐式知识通过正则化图生成模块获取。在每次迭代中，该模块使用Gumbel Softmax对从可训练节点嵌入生成的邻接矩阵进行采样，以类似Dropout的方式丢弃小概率或冗余的边。上述过程可以如下表述：

给定$\boldsymbol{E} \in \mathbb{R}^{N \times d}$为学习的节点嵌入矩阵，$d$为嵌入维度，$\boldsymbol{\theta}$为概率矩阵，$\theta_{ij} \in \boldsymbol{\theta}$表示保留从时间序列$i$到$j$的边的概率，有：
\begin{equation}
\boldsymbol{\theta} = \boldsymbol{E}\boldsymbol{E}^T
\end{equation}

给定$\sigma$为激活函数，$s$为变量，稀疏邻接矩阵$\boldsymbol{A}^{(l)}$定义为：
\begin{equation}
\boldsymbol{A}^{(l)} = \sigma \left( \left( \log\left(\frac{\theta_{ij}}{1-\theta_{ij}}\right) + (g_{ij}^1 - g_{ij}^2) \right) / s \right), \quad g_{ij}^1, g_{ij}^2 \sim \mathrm{Gumbel}(0, 1)
\end{equation}
其中$A_{ij}^{(l)} = 1$对应概率$\theta_{ij}$，0对应剩余概率。

数学上，公式(A-7)使得节点嵌入之间的内积编码能够捕获多变量时间序列之间的隐式相关性，最终产生离散图结构，即隐式图。

\subsubsection*{A.5.2.3\quad 知识融合}

显式知识和隐式知识从不同方面描述了风力发电机组中变量之间的交互，因此二者的适当融合将有利于描述风力发电机组的状态。受文献启发，引入了图操作的切比雪夫多项式展开形式。给定$\boldsymbol{X} \in \mathbb{R}^{N \times H}$为输入矩阵，$\boldsymbol{W}^{(0)}$、$\boldsymbol{W}_b^{(0)}$、$\boldsymbol{W}^{(l)}$和$\boldsymbol{W}_b^{(l)}$为可学习参数，则$\boldsymbol{A}^{(l)}$和$\boldsymbol{A}^{(0)}$的图操作的切比雪夫多项式展开形式为：
\begin{equation}
\tilde{\boldsymbol{A}}^{(0)}(\boldsymbol{X}) = \left(\boldsymbol{I} + \boldsymbol{D}^{-\frac{1}{2}}\boldsymbol{A}^{(0)}\boldsymbol{D}^{-\frac{1}{2}}\right)\boldsymbol{X}\boldsymbol{W}^{(0)} + \boldsymbol{W}_b^{(0)}
\end{equation}
\begin{equation}
\tilde{\boldsymbol{A}}^{(l)}(\boldsymbol{X}) = \left(\boldsymbol{I} + \boldsymbol{D}^{-\frac{1}{2}}\boldsymbol{A}^{(l)}\boldsymbol{D}^{-\frac{1}{2}}\right)\boldsymbol{X}\boldsymbol{W}^{(l)} + \boldsymbol{W}_b^{(l)}
\end{equation}
其中$\boldsymbol{I} \in \mathbb{R}^{N \times N}$，$\boldsymbol{D}$为度矩阵。

受传统Mix-up操作和注意力机制的启发，使用注意力模块度量显式知识和隐式知识之间的相关性，使用Mix-up操作平衡两类知识的组成，以融合显式和隐式知识。

为便于度量这些权重，首先引入读出函数获得隐式和显式知识的嵌入表示$\vec{a}^{(0)}$和$\vec{a}^{(l)}$，连接所有嵌入向量后获得图特征表示矩阵$\boldsymbol{M}$，表示为：
\begin{equation}
\vec{a}^{(0)}(\boldsymbol{X}) = \mathcal{R}\left(\tilde{\boldsymbol{A}}^{(0)}(\boldsymbol{X})\right)
\end{equation}
\begin{equation}
\vec{a}^{(l)}(\boldsymbol{X}) = \mathcal{R}\left(\tilde{\boldsymbol{A}}^{(l)}(\boldsymbol{X})\right)
\end{equation}
\begin{equation}
\boldsymbol{M} = \mathrm{concat}\left(\vec{a}^{(0)}(\boldsymbol{X}), \vec{a}^{(1)}(\boldsymbol{X}), \ldots, \vec{a}^{(l)}(\boldsymbol{X}), \ldots, \vec{a}^{(L)}(\boldsymbol{X})\right)
\end{equation}
其中$\mathcal{R}(\cdot)$为读出函数，此处为Mean函数。$l \in [1, L]$为显式图的编号，$L$为显式图的总数。

不同知识之间的权重系数矩阵$\boldsymbol{S}$由软注意力层计算，公式为：
\begin{equation}
\boldsymbol{S} = \mathrm{softmax}\left(\frac{(\boldsymbol{W}_q\boldsymbol{M})(\boldsymbol{W}_k\boldsymbol{M})^T}{\sqrt{d_k}}\right)
\end{equation}
其中$\mathrm{softmax}(\cdot)$为激活函数，$\boldsymbol{W}_q$和$\boldsymbol{W}_k$为可学习参数，$d_k$为缩放因子。

计算得到的矩阵包含五类注意力权重，含义分别为隐式知识自注意力（记为"$\mathrm{Im} \leftrightarrow \mathrm{Im}$"）、显式知识自注意力（记为"$\mathrm{Ex}_i \leftrightarrow \mathrm{Ex}_i$"）、隐式知识对不同显式知识的注意力（记为"$\mathrm{Im} \rightarrow \mathrm{Ex}_i$"）、不同显式知识对隐式知识的注意力（记为"$\mathrm{Ex}_i \rightarrow \mathrm{Im}$"）以及不同显式知识之间的注意力（记为"$\mathrm{Ex}_i \rightarrow \mathrm{Ex}_j$"（$i \neq j$）），如图~A-4所示。在我们的问题中，隐式知识是在大量数据支持下获得的，在概率意义上具有规律不变性，而显式知识是基于不同显式规则获得的，在实用性意义上具有可变性。因此，本文考虑的知识融合以隐式知识为核心、显式知识为补充来实现。基于上述考虑，三类注意力权重"$\mathrm{Ex}_i \leftrightarrow \mathrm{Ex}_i$"、"$\mathrm{Ex}_i \rightarrow \mathrm{Im}$"和"$\mathrm{Ex}_i \rightarrow \mathrm{Ex}_j$"（$i \neq j$）被忽略，同时"$\mathrm{Im} \leftrightarrow \mathrm{Im}$"的注意力权重被强制设为1，仅以准硬注意力方式关注"$\mathrm{Im} \rightarrow \mathrm{Ex}_i$"的注意力权重。最后，基于这两类过滤后的注意力权重，显式知识$\tilde{\boldsymbol{A}}^{(l)}$和学习到的隐式知识$\tilde{\boldsymbol{A}}^{(0)}$被融合在一起，公式为：
\begin{equation}
\tilde{\boldsymbol{A}}^{(m)}(\boldsymbol{X}) = \tilde{\boldsymbol{A}}^{(0)}(\boldsymbol{X}) + \sum_{l=1}^{L} \alpha_l \tilde{\boldsymbol{A}}^{(l)}(\boldsymbol{X})
\end{equation}
其中$\alpha_l$为隐式知识对不同显式知识的注意力（即"$\mathrm{Im} \rightarrow \mathrm{Ex}_i$"），$\tilde{\boldsymbol{A}}^{(m)}$为融合知识。

% 图A-4: 注意力权重说明
% \begin{figure}[htbp]
% \centering
% \includegraphics[width=0.6\textwidth]{figures/appendixA/fig4_attention.png}
% \caption{注意力权重说明}
% \label{fig:appA-attention}
% \end{figure}

\subsection*{A.5.3\quad 多变量时间序列图神经网络建模与异常程度}

本阶段旨在基于MTGNN构建高精度数字孪生模型。如图~A-5所示，首先使用多变量时间序列数据$\boldsymbol{X}$基于第A.5.2节提出的图构建模块学习融合显式和隐式知识的图邻接矩阵。然后将多变量时间序列数据$\boldsymbol{X}$连同融合知识$\tilde{\boldsymbol{A}}^{(m)}$送入MTGNN。在通过MTGNN的图卷积模块和时序卷积模块分别提取数据中的变量间依赖关系和时序依赖关系后，输出模块给出下一时刻的预测值$\hat{\boldsymbol{Y}}$。最后，使用绝对损失优化目标值$\boldsymbol{Y}$与预测值$\hat{\boldsymbol{Y}}$之间的差异，以构建高精度模型。该训练过程可以表述为：
\begin{equation}
\hat{\boldsymbol{Y}} = f\left(\boldsymbol{X}, \tilde{\boldsymbol{A}}^{(m)}; \boldsymbol{\Theta}\right)
\end{equation}
\begin{equation}
L = \mathrm{Loss}(\hat{\boldsymbol{Y}}, \boldsymbol{Y})
\end{equation}
其中$f(\cdot)$为从$\boldsymbol{X}$到$\boldsymbol{Y}$的映射函数，$\boldsymbol{\Theta}$为$f(\cdot)$的可学习参数，$\mathrm{Loss}(\cdot)$为绝对损失函数。

% 图A-5: MTGNN建模的概念图
% \begin{figure}[htbp]
% \centering
% \includegraphics[width=0.8\textwidth]{figures/appendixA/fig5_concept.png}
% \caption{MTGNN建模的概念图}
% \label{fig:appA-concept}
% \end{figure}

异常程度的定义来源于正常工况下模型预测偏差的分布。此处假设该分布为高斯分布。根据高斯分布的参数（即均值和标准差），定义变量状态的异常程度，如图~A-6所示。

% 图A-6: 异常程度的定义
% \begin{figure}[htbp]
% \centering
% \includegraphics[width=0.7\textwidth]{figures/appendixA/fig6_anomaly.png}
% \caption{异常程度的定义}
% \label{fig:appA-anomaly}
% \end{figure}

\subsection*{A.5.4\quad 根因定位与故障传播分析}

根因定位是对故障根源的探索，依赖于故障传播分析。通过可视化节点状态随时间的变化，各监测变量的状态变化将被明确，故障状态在不同节点（监测变量）之间的传播也将被确定。最后，通过在反向时间轴上的回溯分析实现根因定位，如图~A-7所示。

% 图A-7: 故障传播分析与根因定位
% \begin{figure}[htbp]
% \centering
% \includegraphics[width=0.7\textwidth]{figures/appendixA/fig7_propagation.png}
% \caption{故障传播分析与根因定位}
% \label{fig:appA-propagation}
% \end{figure}

\section*{A.6\quad 案例研究}
\addcontentsline{toc}{section}{A.6\quad 案例研究}

\subsection*{A.6.1\quad 数据集与实验设置}

案例研究使用了华北某风电场采集的SCADA数据。该SCADA系统监测了一台风力发电机组的50个运行参数，包括1秒平均风速、有功功率、发电机平均转速、齿轮箱轴承温度和齿轮箱油温等，采样间隔为10分钟。该数据集的时间跨度为2019年1月至2020年12月。本研究选取了23个监测参数用于后续研究，其名称和符号表示如表~A-1所示。在数据预处理阶段，首先基于标准风速功率曲线对SCADA数据集进行清洗，然后进行最大值归一化处理。数据集按0.6、0.2、0.2的比例划分为训练集、验证集和测试集。

\begin{table}[htbp]
\centering
\caption{风力发电机组23个参数及其符号表示}
\label{tab:appA-params23}
\small
\begin{tabular}{clccl c}
\toprule
编号 & 参数 & 单位 & 编号 & 参数 & 单位 \\
\midrule
M1 & NCC300温度 & $^\circ$C & M13 & 最大发电机绕组温度 & $^\circ$C \\
M2 & NCC320温度 & $^\circ$C & M14 & 变桨电机1温度 & $^\circ$C \\
M3 & 机侧电感温度 & $^\circ$C & M15 & 变桨电机2温度 & $^\circ$C \\
M4 & 网侧电感温度 & $^\circ$C & M16 & 变桨电机3温度 & $^\circ$C \\
M5 & 机侧半导体温度 & $^\circ$C & M17 & 偏航变频器温度 & $^\circ$C \\
M6 & 网侧半导体温度 & $^\circ$C & M18 & 机舱温度 & $^\circ$C \\
M7 & 环境温度 & $^\circ$C & M19 & 有功功率 & kW \\
M8 & 滤板温度 & $^\circ$C & M20 & 机舱电池温度 & $^\circ$C \\
M9 & 齿轮箱轴承温度 & $^\circ$C & M21 & 1秒平均风速 & m/s \\
M10 & 齿轮箱油温 & $^\circ$C & M22 & 发电机平均转速 & rpm \\
M11 & 驱动方向发电机轴承温度 & $^\circ$C & M23 & 控制发电机平均转速 & rpm \\
M12 & 非驱动方向发电机轴承温度 & $^\circ$C & & & \\
\bottomrule
\end{tabular}
\end{table}

为展示GNN的优势及所提出的显式-隐式知识融合模块带来的性能提升，选取了10种基准方法进行建模精度比较，包括LSTM、BiLSTM、CNN-LSTM、Transformer、GCRN、STGCN、GraphLSTM、GAT、NRI和本文使用的MTGNN。MTGNN已在第A.4节详细描述，其他九种方法简述如下：

\textbf{LSTM：}长短期记忆网络是循环神经网络的变体，可以通过门控将短期记忆与长期记忆相结合，解决传统循环神经网络中的梯度消失和梯度爆炸问题。

\textbf{BiLSTM：}双向长短期记忆网络是LSTM的变体，结合前向LSTM和后向LSTM以更好地处理输入特征。

\textbf{CNN-LSTM：}通过简单堆叠1D-CNN和LSTM获得的混合模型。

\textbf{Transformer：}一种完全依赖自注意力机制来计算输入和输出表示的神经网络。

\textbf{GCRN：}图卷积循环网络，作为传统循环神经网络在任意结构上的广义形式，是一种能够预测结构化序列数据的深度学习模型。

\textbf{STGCN：}时空图卷积网络，是由时空卷积块构建的经典时空图模型，可以自动学习数据中的空间和时序特征。

\textbf{GraphLSTM：}将LSTM从时序或多维数据推广到图结构数据的神经网络。

\textbf{GAT：}图注意力网络，是基于注意力机制的图神经网络。

\textbf{NRI：}神经关系推理模型，是一种能够以无监督方式学习交互系统中动态模型以推断隐藏交互结构的神经网络。

使用均方根误差（RMSE）和决定系数（$R^2$）两个指标评价建模精度，RMSE越小、$R^2$越大，预测精度越高：
\begin{equation}
\mathrm{RMSE} = \sqrt{\frac{1}{n}\sum_{i=1}^{n}(\hat{y}_i - y_i)^2}
\end{equation}
\begin{equation}
R^2 = 1 - \frac{\sum_{i=1}^{n}(\hat{y}_i - y_i)^2}{\sum_{i=1}^{n}(\bar{y} - y_i)^2}
\end{equation}
其中$n$为样本数量，$\hat{y}_i$为$i$时刻的预测值，$y_i$为同一时刻的监测值，$\bar{y}$为样本均值。

\subsection*{A.6.2\quad 结果}

\subsubsection*{A.6.2.1\quad 图的研究}

\textbf{i) 显式图}

风力发电机组中的机械传动关系（$G_1$）涉及风速、发电机转速和有功功率等变量，用于描述风力发电机组捕获风能后依次将其转化为机械能和电能的过程。有向边的连接如表~A-2所示。

\begin{table}[htbp]
\centering
\caption{由机械传动关系定义的边}
\label{tab:appA-mechanical}
\small
\begin{tabular}{ll}
\toprule
边的类型 & 边~$(u, v)$ \\
\midrule
风能转化为机械能 & M21 $\rightarrow$ M22，M21 $\rightarrow$ M23 \\
机械能转化为电能 & M22 $\rightarrow$ M19，M23 $\rightarrow$ M19 \\
\bottomrule
\end{tabular}
\end{table}

风力发电机组中的温度影响关系（$G_2$）涉及许多温度相关变量，用于描述变量之间的温度交互。有向边的连接如表~A-3所示。

\begin{table}[htbp]
\centering
\caption{由温度影响关系定义的边}
\label{tab:appA-temperature}
\small
\begin{tabular}{lp{9cm}}
\toprule
边的类型 & 边~$(u, v)$ \\
\midrule
环境温度对变桨系统和机舱温度的影响 & M7 $\rightarrow$ M14，M7 $\rightarrow$ M15，M7 $\rightarrow$ M16，M7 $\rightarrow$ M18 \\
\addlinespace
机舱温度与内部组件温度的交互 & M18 $\leftrightarrow$ M1，M18 $\leftrightarrow$ M2，M18 $\leftrightarrow$ M3，M18 $\leftrightarrow$ M4，M18 $\leftrightarrow$ M5，M18 $\leftrightarrow$ M6，M18 $\leftrightarrow$ M8，M18 $\leftrightarrow$ M9，M18 $\leftrightarrow$ M10，M18 $\leftrightarrow$ M11，M18 $\leftrightarrow$ M12，M18 $\leftrightarrow$ M13，M18 $\leftrightarrow$ M20，M18 $\leftrightarrow$ M17 \\
\addlinespace
传动链上的温度影响 & M21 $\rightarrow$ M14，M21 $\rightarrow$ M15，M21 $\rightarrow$ M16，M21 $\rightarrow$ M9，M21 $\rightarrow$ M10，M9 $\rightarrow$ M11，M10 $\rightarrow$ M11，M22 $\rightarrow$ M11，M22 $\rightarrow$ M12，M22 $\rightarrow$ M13 \\
\addlinespace
NCC300控制柜中的温度影响 & M1 $\rightarrow$ M17 \\
\addlinespace
NCC320控制柜中的温度影响 & M2 $\leftrightarrow$ M3，M2 $\leftrightarrow$ M4，M2 $\leftrightarrow$ M5，M2 $\leftrightarrow$ M6，M2 $\leftrightarrow$ M8，M3 $\leftrightarrow$ M4，M3 $\leftrightarrow$ M5，M3 $\leftrightarrow$ M6，M3 $\leftrightarrow$ M8，M4 $\leftrightarrow$ M5，M4 $\leftrightarrow$ M6，M4 $\leftrightarrow$ M8，M5 $\leftrightarrow$ M6，M5 $\leftrightarrow$ M8，M6 $\leftrightarrow$ M8 \\
\bottomrule
\end{tabular}
\end{table}

基于相关性（$G_3$）、互信息（$G_4$）和因果性（$G_5$）的图通过第A.5.2.1节中说明的方法获得。

% 图A-9: 显式图与融合图
% \begin{figure}[htbp]
% \centering
% \includegraphics[width=\textwidth]{figures/appendixA/fig9_graphs.png}
% \caption{显式图与融合图}
% \label{fig:appA-graphs}
% \end{figure}

\textbf{ii) 融合图}

使用第A.5.2.3节介绍的算法，以隐式知识为核心、不同权重的显式知识为补充，生成了融合图（$G_{fused}$）。与所提出的五个显式图相比，融合图包含更丰富的信息，这反映了显式知识的局限性和不完整性，换言之，数据中还有许多隐式边信息等待被挖掘。

值得注意的是，融合图中的一些边信息是可解释的。例如，"M19 $\rightarrow$ M21"的边表达了风速与功率之间的关系，该关系未在显式知识中直接定义，但在挖掘隐式知识后从数据中成功学习得到。"M9 $\leftrightarrow$ M10"的边表达了齿轮箱轴承温度与齿轮箱油温之间的交互，同样未直接定义，但仍被所提算法学习到。这些典型关系的学习为知识融合的有效性提供了有力证据，并支持了后续风力发电机组数字孪生模型的构建。

\subsubsection*{A.6.2.2\quad 建模精度}

使用第A.6.1节列出的两个评价指标比较了多种基准模型的建模精度，不同模型的输入输出参数配置相同。不同模型（传统时间序列神经网络和时间序列GNN）在无图、五个显式图和融合图七种条件下的时间序列预测结果如表~A-4所示。

\begin{table}[htbp]
\centering
\caption{10种模型在2个评价指标（RMSE和$R^2$）下的建模精度}
\label{tab:appA-accuracy}
\small
\begin{tabular}{lcccccccc}
\toprule
\multirow{2}{*}{模型} & \multicolumn{7}{c}{RMSE} \\
\cmidrule{2-8}
 & 无图 & $G_1$ & $G_2$ & $G_3$ & $G_4$ & $G_5$ & $G_{fused}$ \\
\midrule
LSTM & 0.0731 & --- & --- & --- & --- & --- & --- \\
BiLSTM & 0.0752 & --- & --- & --- & --- & --- & --- \\
CNN-LSTM & 0.0772 & --- & --- & --- & --- & --- & --- \\
Transformer & 0.0699 & --- & --- & --- & --- & --- & --- \\
GraphLSTM & --- & 0.0631 & 0.0641 & 0.0632 & 0.0636 & 0.0635 & 0.0629 \\
GAT & --- & 0.0970 & 0.1194 & 0.1556 & 0.0899 & 0.0964 & 0.1026 \\
NRI & --- & 0.0417 & 0.0416 & 0.0378 & 0.0416 & 0.0417 & 0.0416 \\
GCRN & --- & 0.0443 & 0.0479 & 0.0463 & 0.0496 & 0.0460 & 0.0459 \\
STGCN & --- & 0.0697 & 0.0591 & 0.0624 & 0.0679 & 0.0694 & 0.0600 \\
MTGNN & 0.0525 & 0.0519 & 0.0530 & 0.0541 & 0.0516 & 0.0438 & \textbf{0.0427} \\
\bottomrule
\end{tabular}

\vspace{6pt}

\begin{tabular}{lcccccccc}
\toprule
\multirow{2}{*}{模型} & \multicolumn{7}{c}{$R^2$} \\
\cmidrule{2-8}
 & 无图 & $G_1$ & $G_2$ & $G_3$ & $G_4$ & $G_5$ & $G_{fused}$ \\
\midrule
LSTM & 0.8346 & --- & --- & --- & --- & --- & --- \\
BiLSTM & 0.8485 & --- & --- & --- & --- & --- & --- \\
CNN-LSTM & 0.8330 & --- & --- & --- & --- & --- & --- \\
Transformer & 0.8224 & --- & --- & --- & --- & --- & --- \\
GraphLSTM & --- & 0.9192 & 0.9143 & 0.9096 & 0.9163 & 0.9164 & 0.9145 \\
GAT & --- & 0.7464 & 0.5942 & 0.3547 & 0.7839 & 0.7378 & 0.7271 \\
NRI & --- & 0.8569 & 0.9623 & 0.9469 & 0.8825 & 0.9605 & 0.9631 \\
GCRN & --- & 0.9647 & 0.9517 & 0.9459 & 0.9518 & 0.9557 & 0.9305 \\
STGCN & --- & 0.8687 & 0.9188 & 0.9016 & 0.8793 & 0.8781 & 0.9136 \\
MTGNN & 0.9440 & 0.9440 & 0.9428 & 0.9406 & 0.9458 & 0.9655 & \textbf{0.9666} \\
\bottomrule
\end{tabular}
\end{table}

可以看出，使用图和时间序列数据构建的时间序列GNN比仅使用时间序列数据构建的时间序列神经网络具有更高的精度。这可以解释为时间序列GNN以更显式的形式将变量间交互纳入训练过程，网络更有可能在正确的方向上被训练；而时间序列神经网络以隐式方式考虑变量间交互，完全参数化的交互学习具有更大的自由度且更难学习，使得建模精度难以超过时间序列GNN。

此外，融合知识相比显式知识对建模精度的提升也是关注重点。如表~A-4所示，基于融合图构建的模型具有最高（或接近最高）的精度（即最小RMSE和最大$R^2$），这意味着融合图帮助这些模型更准确地描述了风力发电机组的正常状态。这可以解释为显式图只能为模型训练提供部分有效信息，而融合图在提供变量间交互信息方面更加全面。

\subsubsection*{A.6.2.3\quad 故障传播分析与根因定位}

本节使用实际风电场的三个故障案例验证所提框架的有效性：i) 有功功率偏差大于或等于500~kW；ii) 三个控制柜中两个温度过高；iii) 滤板温度超过100度。

\textbf{i) 有功功率偏差大于或等于500~kW}

该故障与参数M19"有功功率"相关。当告警发出时，参数M19的异常程度已达到最大值1，因此该故障的告警是合理的。回溯故障传播过程可以发现，在故障发生前的半小时内，仅有参数M11"驱动方向发电机轴承温度"的异常程度在逐渐增加，直到故障发生时刻，M11引起的异常沿系统网络扩展到其他参数，最终导致M19的异常告警。因此，可以将该案例中故障的根因定位为发电机轴承异常。

\textbf{ii) 三个控制柜中两个温度过高}

该故障与参数M1"NCC300温度"或M2"NCC320温度"相关，因此需要进一步确认故障的源节点。当故障告警发出时，M2的异常程度已达到最大值1，与故障描述一致。回溯故障传播过程可以看出，在故障前一小时内，仅有M2的异常程度在逐渐增加，然后M2引起的异常扩展到其他参数并最终导致告警。因此，可以将该故障的根因定位为控制柜"NCC320"的温度异常。

\textbf{iii) 滤板温度超过100度}

该故障与参数M8"滤板温度"相关。当故障告警发出时，M8的异常程度已达到最大值1，与故障描述一致。回溯故障传播过程可以看出，M1"NCC300温度"引起的异常首先扩展到M18"机舱温度"，然后扩展到M8，最终触发告警。因此，可以将该故障的根因定位为控制柜"NCC300"的温度异常。

% 图A-11至A-13为三个故障案例的传播分析图
% 受限于篇幅，此处省略故障传播图的详细展示

\subsection*{A.6.3\quad 卢布林糖厂、磨煤机与另一台风力发电机组的案例研究}

所提框架基于预测偏差的异常检测策略因其简洁性和可靠性而具有高度泛化能力。为证明这一点，本文使用卢布林糖厂公开数据集、某火力发电厂的磨煤机监测数据集以及另一个风电场的风力发电机组SCADA数据集的故障样本作为案例研究。为降低建模难度并进一步提高框架的泛化能力，此处仅将相关性、互信息和因果性作为显式知识与隐式知识融合，然后为卢布林糖厂系统、磨煤机或风力发电机组构建高精度数字孪生模型，以进行特定案例的根因定位。

\subsubsection*{A.6.3.1\quad 卢布林糖厂案例}

DAMADICS基准提供的卢布林糖厂数据由32个测量变量组成，采样率为1~Hz，时间跨度为2001年10月29日至11月22日。32个参数的符号表示和名称如表~A-5所示。

\begin{table}[htbp]
\centering
\caption{卢布林糖厂32个参数及其符号表示}
\label{tab:appA-sugar}
\small
\begin{tabular}{clccl c}
\toprule
编号 & 参数 & 单位 & 编号 & 参数 & 单位 \\
\midrule
N1 & P1--果汁压力(阀入口) & kPa & N17 & P1--果汁压力(阀入口) & kPa \\
N2 & P2--果汁压力(阀出口) & kPa & N18 & P2--果汁压力(阀出口) & kPa \\
N3 & T--果汁温度(阀入口) & $^\circ$C & N19 & T--果汁温度(阀入口) & $^\circ$C \\
N4 & F--果汁流量(1st蒸发器入口) & m$^3$/h & N20 & PV--过程值(果汁流量，5th蒸发器出口) & m$^3$/h \\
N5 & CV--控制值(控制器输出) & \% & N21 & CV--控制值(控制器输出) & \% \\
N6 & X--伺服电机杆位移 & \% & N22 & X--伺服电机杆位移 & \% \\
N7 & PV--过程值(1st蒸发器液位) & \% & N23 & P1--水压力(阀入口) & kPa \\
N8 & 果汁温度(1st蒸发器入口) & $^\circ$C & N24 & P2--水压力(阀出口) & kPa \\
N9 & 果汁温度(1st蒸发器出口) & $^\circ$C & N25 & T--水温(阀出口) & $^\circ$C \\
N10 & 果汁密度(1st蒸发器入口) & BX & N26 & F--水流量(蒸汽锅炉入口) & t/h \\
N11 & 果汁密度(1st蒸发器出口) & BX & N27 & CV--控制值(控制器输出) & \% \\
N12 & 蒸汽流量 & t/h & N28 & X--伺服电机杆位移 & \% \\
N13 & 蒸汽压力 & kPa & N29 & PV--过程值(蒸汽锅炉水位) & \% \\
N14 & 蒸汽温度 & $^\circ$C & N30 & 蒸汽流量(蒸汽锅炉出口) & t/h \\
N15 & 蒸气压力 & kPa & N31 & 蒸汽压力(蒸汽锅炉出口) & kPa \\
N16 & 蒸气温度 & $^\circ$C & N32 & 蒸汽温度(蒸汽锅炉出口) & $^\circ$C \\
\bottomrule
\end{tabular}
\end{table}

人工故障"执行器3的定位器供给压力降"用于案例研究。根据数据描述文件，该人工故障于2001年11月17日引入，表现为参数N28"X--伺服电机杆位移"在故障前20秒开始下降，偏离正常运行工况。基于所提框架，成功评估了故障前糖厂系统中参数状态的传播过程。可以看出，参数N28的状态在故障前16秒变得明显异常，然后该异常传播到其他节点，最终在故障时刻达到最大异常程度。在根因定位方面，回溯上述故障传播过程，可以判断该糖厂系统人工故障的根因为参数N28异常，这与事实一致，证明了该框架能够实现根因定位。

\subsubsection*{A.6.3.2\quad 磨煤机案例}

华东某火力发电厂提供的磨煤机数据由30个测量变量组成，采样率为1/30~Hz，时间跨度为2021年7月至2022年3月。30个参数的符号表示和名称如表~A-6所示。

\begin{table}[htbp]
\centering
\caption{磨煤机30个参数及其符号表示}
\label{tab:appA-coal}
\small
\begin{tabular}{clcclc}
\toprule
编号 & 参数 & 单位 & 编号 & 参数 & 单位 \\
\midrule
V1 & 磨煤机电流 & A & V16 & 磨煤机电机B相定子绕组温度2 & $^\circ$C \\
V2 & 磨煤机润滑油泵电流 & A & V17 & 磨煤机电机C相定子绕组温度1 & $^\circ$C \\
V3 & 磨煤机出口压力1 & kPa & V18 & 磨煤机电机C相定子绕组温度2 & $^\circ$C \\
V4 & 磨煤机出口压力2 & kPa & V19 & 磨煤机电机轴端轴承温度 & $^\circ$C \\
V5 & 磨煤机出口温度1 & $^\circ$C & V20 & 磨煤机电机非轴端轴承温度 & $^\circ$C \\
V6 & 磨煤机出口温度2 & $^\circ$C & V21 & 磨煤机碾磨碗上/下差压 & kPa \\
V7 & 磨煤机行星齿轮箱轴承温度1 & $^\circ$C & V22 & 磨煤机密封风与一次风差压 & \% \\
V8 & 磨煤机行星齿轮箱轴承温度2 & $^\circ$C & V23 & 磨煤机冷风挡板开度反馈 & \% \\
V9 & 磨煤机行星齿轮箱轴承温度3 & $^\circ$C & V24 & 磨煤机热风挡板开度反馈 & \% \\
V10 & 磨煤机行星齿轮箱轴承温度4 & $^\circ$C & V25 & 磨煤机一次风流量 & km$^3$/h \\
V11 & 磨煤机润滑油箱温度 & $^\circ$C & V26 & 磨煤机一次风流量设定值 & km$^3$/h \\
V12 & 磨煤机润滑油温度 & $^\circ$C & V27 & 磨煤机出口温度设定值 & $^\circ$C \\
V13 & 磨煤机电机A相定子绕组温度1 & $^\circ$C & V28 & 磨煤机入口一次风流量 & km$^3$/h \\
V14 & 磨煤机电机A相定子绕组温度2 & $^\circ$C & V29 & 磨煤机入口一次风温度 & $^\circ$C \\
V15 & 磨煤机电机B相定子绕组温度1 & $^\circ$C & V30 & 磨煤机一次风流量 & km$^3$/h \\
\bottomrule
\end{tabular}
\end{table}

实际故障"磨煤机磨辊轴承故障"用于案例研究。根据磨煤机结构，该故障会导致磨煤机齿轮箱和电机的温度测量异常，如V7--V10、V13--V20等。基于所提框架，成功评估了故障前磨煤机中参数状态的传播过程。可以看出，在故障时刻参数V8、V10、V13--V18、V20和V29表现出较高的异常程度，与事实一致。回溯故障传播过程后，上述参数均反映了异常程度随时间逐渐增加的趋势，其中参数V20"磨煤机电机非轴端轴承温度"的状态最早变得异常。因此，可以判断该故障的根因为参数V20异常，实际检修中应优先关注与V20相关的组件。

\subsubsection*{A.6.3.3\quad 另一个风电场的风力发电机组案例}

该案例由中国东北某风电场（不同于前述风电场）提供的监测数据组成，包含50个测量变量，采样间隔为10分钟，时间跨度为2019年1月至2021年12月。

实际故障"变频器异常"用于案例研究。根据测量传感器的位置，该故障可与变频器控制柜中的监测参数相关，如M1--M6、M9、M10等。基于所提框架，成功评估了该故障前风力发电机组中参数状态的传播过程。当监测系统告警时，M5、M6、M9、M10、M19--M21、M25、M31、M39--M41、M45--M50等参数表现出较高的异常程度，而在故障前的时刻，仅有参数M9、M31、M42--M44、M46和M49是异常的。因此，该故障的根因可以定位在叶片、变频器柜或发电机，可在后续维护过程中重点关注。

\subsection*{A.6.4\quad 讨论}

\textbf{物理引导的显式知识。}对于GNN的实际应用，"使用预定义图进行模型训练"是一项困难的任务，因为在许多工业场景中（如为设备构建数字孪生模型）预定义图并不存在。因此，如何使用某些规则显式地获得预定义图是一个重要关切。对于风力发电机组数字孪生模型构建的场景，本文提供了一种通过从五个方面定义描述变量间交互的系列显式知识来产生预定义图的思路，包括机械传动关系、温度影响关系、相关性、互信息和因果性，这同时也为在构建设备数字孪生模型的工业场景中获得预定义图提供了解决方案。

\textbf{数据驱动的隐式知识与知识融合。}需要注意的是，显式知识不能完全正确地表达图中节点之间的关系。为获得更大概率匹配真实世界场景的预定义图，需要从数据中提取隐式知识并融合这两类知识。为实现这一目标，本文基于前期研究，以隐式知识为核心、以不同权重的显式知识为补充，以准硬注意力方式实现知识融合。该方法可以嵌入到不同场景的显式-隐式知识融合过程中，有效获得GNN训练所需的预定义图。

\textbf{时空特征学习与建模精度。}在本文中，风力发电机组的数字孪生模型构建被抽象为多变量时间序列预测。相比于传统时间序列神经网络更多关注时序依赖关系而不考虑或仅隐式考虑变量间依赖关系的局限性，同时考虑变量间依赖关系和时序依赖关系的时间序列图神经网络具有优越的预测性能。因此，本文提出了一种使用从知识融合模块获得的显式-隐式融合图和MTGNN的高精度数字孪生模型构建框架，其建模精度（RMSE：0.0427，$R^2$：0.9666）与传统时间序列神经网络和其他时间序列图神经网络进行了比较，一方面验证了时间序列图神经网络在时间序列预测方面的优势，另一方面验证了融合的显式-隐式知识对时间序列图神经网络预测性能的提升。实际上，该框架具有很强的泛化能力，不同设备/系统（如本文中的磨煤机和卢布林糖厂）的高精度数字孪生模型均可用同样的思路来构建。

\textbf{故障传播分析与根因定位。}在本文中，故障传播分析首先通过准确描述每个参数在每个时刻的异常程度（偏离正常状态的程度），然后逐时刻分析故障状态在参数之间的传播来实现。故障传播的研究在学术和实践方面都具有重要价值。一方面，该研究通过调查故障模式在复杂系统中的传播，使故障诊断从"是什么"的研究转变为"为什么"的研究，并进一步有助于发现故障传播机制。另一方面，当该研究与工业实践相结合时，工厂操作人员将同时关注故障设备（当前正在告警的设备）和故障源设备（最初产生故障模式的设备），这将提高计划内系统化运维水平并降低非计划停机风险。

\textbf{低采样率数据上的应用与框架泛化能力。}现实世界的各种场景中存在大量采样率低、正常工况占比高、成本低的多变量监测数据，如何有效利用这些数据进行设备故障预警、检测、定位和诊断是一个备受关注且困难的问题。为此，本文提出了一种处理此类数据并将其用于设备运维的有效解决方案，即在正常工况下构建高精度数字孪生模型，通过跟踪实际设备状态相对于正常状态的偏差来实现异常检测和根因定位。在案例研究中，本文成功复现了风力发电机组、磨煤机和糖厂系统中实际故障的传播过程，并通过回溯分析实现了根因定位。总之，本研究不仅提供了如何利用低采样率监测数据实现故障定位的思路，而且说明了基于预测偏差的异常检测策略以其简洁性和可靠性可以确保所提框架的泛化能力。

\section*{A.7\quad 结论}
\addcontentsline{toc}{section}{A.7\quad 结论}

受利用大量监测数据实现风力发电机组根因定位的迫切需求驱动，本文充分利用GSL强大的变量间依赖关系挖掘能力和时间序列图神经网络卓越的多变量时间序列预测能力，提出了一种基于显式-隐式知识融合和MTGNN的风力发电机组根因定位框架。基于五种显式规则（包括机械传动关系、温度影响关系、相关性、互信息和因果性）构建了风力发电机组变量间的显式知识，以显式方式表达变量间依赖关系。基于GSL思想，从数据中学习到的隐式知识与显式知识以准硬注意力方式进行加权融合，以获得GNN所需的图邻接矩阵。利用MTGNN框架，基于融合知识训练了高精度数字孪生模型。最后，定义了节点状态的异常程度以准确刻画故障状态随时间的传播，此后通过回溯分析实现根因定位。利用实际风力发电机组、磨煤机和糖厂的故障案例验证了所提框架，结果表明：

i) 由五种显式规则定义的显式知识可以部分描述风力发电机组中的变量间依赖关系。

ii) 通过GSL算法和使用准硬注意力机制的加权融合模块学习到的融合图可以更丰富地表达变量间依赖关系，并提高MTGNN的建模精度。

iii) 基于高精度风力发电机组数字孪生模型和定义的节点状态异常程度，所提框架可以有效描述故障状态的传播过程，并通过回溯分析成功实现根因定位。

iv) 所提出的根因定位框架可以广泛应用于不同设备以提高运维效率并降低成本。

在未来的研究中，可以使用不同的预警算法分析多变量预测偏差，以增强数字孪生模型的预警能力。
